\BourbakiExerciseGroup

\begin{enumerate}
\def\labelenumi{\arabic{enumi})}
\item
  设 \(\mathcal{C}_1, \mathcal{C}_2\) 是同一集合 X 上的两个拓扑。证明：\(\mathcal{C}_1\) 比 \(\mathcal{C}_2\) 更细，当且仅当 X 上每个在拓扑 \(\mathcal{C}_1\) 中收敛的滤子都在拓扑 \(\mathcal{C}_2\) 中收敛到同一点。
\item
  设 \(\mathfrak{U}\) 是 \(\mathbf{N}\) 上比 Fréchet 滤子更细的一个超滤子，并设 \(\mathbf{X} = \mathbf{N} \cup \{\omega\}\) 为与 \(\mathfrak{U}\) 相伴随的空间（§ 6，第 5 目）。证明：具有无穷多个不同项的序列 \((x_n)\) 在 \(\mathbf{X}\) 中决不收敛。
\item
  设 X, \(X'\) 为两个拓扑空间，\(f: X \to X'\) 为在点 \(x_{0} \in X\) 处连续的映射。证明：若 B 是 X 上以 \(x_{0}\) 为聚点的任一滤子基，则 \(f(x_{0})\) 是滤子基 \(f(\mathfrak{B})\) 在 \(X'\) 上的聚点。
\item
  \begin{enumerate}
  \def\labelenumii{\alph{enumii})}
  \tightlist
  \item
    设 X, Y 为两个拓扑空间，F 为 X 上以 a 为聚点的滤子，G 为 Y 上以 b 为聚点的滤子。证明 \((a, b)\) 是积滤子 \(F \times G\) 的聚点。
  \end{enumerate}
\end{enumerate}

\begin{enumerate}
\def\labelenumi{\alph{enumi})}
\setcounter{enumi}{1}
\tightlist
\item
  在积空间 \(Q^{2}\) 中，给出一个序列 \((x_{n}, y_{n})\) 的例子：它没有聚点，尽管序列 \((x_{n})\) 与 \((y_{n})\) 各自在 Q 中都有聚点。
\end{enumerate}

\begin{enumerate}
\def\labelenumi{\arabic{enumi})}
\setcounter{enumi}{4}
\item
  设 X, Y 为两个拓扑空间，G 是 \(X \times Y\) 中映射 \(f: X \to Y\) 的图像。证明：对每个 \(x \in X\)，f 在 x 处的聚点所成的集就是 \(\overline{\mathbf{G}}(x)\)，即 \(\overline{\mathbf{G}}\)（G 在 \(X \times Y\) 中的闭包）在 x 处的截口。
\item
  设 \((\mathbf{X}_t)_{t\in I}\) 是由离散空间组成的一个不可数族，其中每个空间至少有两个点。在积集合上考虑拓扑
\end{enumerate}

\[
\mathbf {X} = \prod_ {\iota \in \mathbf {I}} \mathbf {X} _ {\iota},
\]

它由所有的积 \(\prod_{i\in I}M_i\) 生成，其中 \(M_i\subset X_i\) 对一切 \(i\in I\) 成立，且除去可数个指标外均有 \(M_i = X_i\)（参见§ 4，习题 9）。证明：开集的每个可数交在 \(X\) 中都是开集，并由此推出，\(x\) 的任何点都不具有可数的基本邻域系。

\begin{enumerate}
\def\labelenumi{\arabic{enumi})}
\setcounter{enumi}{6}
\tightlist
\item
  拓扑空间 X 的子集 A 称为本原的，如果它是 X 上某个超滤子的极限点集。设 \(Y_{A}\) 表示 X 上以 A 为极限点集的超滤子所成的集。
\end{enumerate}

\begin{enumerate}
\def\labelenumi{\alph{enumi})}
\item
  若 A 是 X 的本原子集，证明 X 的每个与 A 相交的开子集都属于所有超滤子 \(U \in Y_{A}\)。
\item
  设 A, B 是 X 的两个不同的本原子集。证明存在超滤子 \(U \in Y_{A}\)、超滤子 \(B \in Y_{B}\) 以及 X 的子集 M，使得 \(M \in U\) 且 \(CM \in B\)。
\item
  设 \(f: \mathbf{X} \to \mathbf{Y}\) 是连续映射。证明在 \(f\) 下，\(\mathbf{X}\) 的任何本原子集的像都含于 \(\mathbf{Y}\) 的某个本原子集。
\item
  若 \(x \in \mathbf{X}\) 且 \(\{x\}\) 在 \(\mathbf{X}\) 中闭，证明 \(\{x\}\) 是 \(\mathbf{X}\) 的本原子集。
\end{enumerate}

